\section{国际海上货物运输法律制度}
\begin{enumerate}
  \item 货物运输法要求的多一些。海运法，国际铁路货物运输法，国际航空货物运输法。
  \item 运输方式。
        \begin{enumerate}
          \item 班轮运输。类似地铁。固定的始发港目的港。小中货间杂货。
          \item 租船运输。打出租。航次跑一个航程，定期跑一段时间。租船合同。
        \end{enumerate}
  \item 单证类型。
        \begin{kk}{海运提单}
          用以证明海上货物运输合同的订立和货物已经由承运人接管 \hz 装船以及承运人保证据以交付货物的单证。
          \begin{enumerate}
            \item 是\zd{海上运输合同的证明}。
            \item 是承运人出具的接收货物的收据。
            \item 是承运人交付货物的\zd{物权凭证}。见正本提单放货。\qd{转让海运提单意味着转让在途运输的货物所有权}。
          \end{enumerate}
          \tcblower
          分类标准。
          \begin{enumerate}
            \item 是否已实际装船。已装船提单（船名和航次），收货待运提单/备运提单（船名航次空着）。
            \item 有无货物不良批注。清洁提单和不清洁提单。如少一箱，部分纸箱破损。
            \item 收货人抬头不同。记名/不记名/指示提单。不能转让/交付即转让/背书方能转让。
            \item 出倒签提单和预借提单。把装船时间往前倒着签。装船时间假的。用的少。考过一次。
          \end{enumerate}
        \end{kk}

        \begin{kk}{海运单}
          航程较短的运输。单货一起。不是物权凭证。不能单方转让。都具有运输合同证明和货物收据的作用。
        \end{kk}
  \item 承运人责任。\qd{违反最低限度的义务 \yjt 无过失 \yjt 航行过失 \yjt 火灾过失 \yjt 看规则}。
        \begin{enumerate}
          \item \zd{交货责任}。货交给谁，交错了怎么办，赔不赔怎么赔。
                \begin{enumerate}
                  \item \zd{无单放货}。\zd{无正本提单}交付货物案件。
                  \item 与无单取货人承担\zd{连带责任}。
                  \item 告承运人。违约和侵权竞合。
                  \item 赔偿额以货物的 CIF 价格计算。足额赔偿（\zd{防止恶意串通}）。
                  \item 免责。
                        \begin{enumerate}
                          \item 海关、法院、商检机关等依法处理，过时效无人清关。依法行为/依法免责。
                          \item 执行\zd{记名提单}托运人指示（终止运输、返还货物、变更到达地、交其他收货人）。
                          \item 已有正本提单已提货。一式数份正本提单（会计入账需求）。只认第一份。
                        \end{enumerate}
                  \item 诉讼时效 \zd{1} 年，适用法律为中国海商法。
                \end{enumerate}
          \item \zd{货损责任}。在运输途中出现了毁损灭失，赔不赔，怎么赔。海牙规则、维斯比规则、汉堡规则。
                \begin{kk}{海牙规则、维斯比规则和汉堡规则}
                  \begin{enumerate}
                    \item 承运人最低限度的义务。适航义务，管货义务。\zd{所有海运规则都包含}。
                    \item 责任期间。装上船到卸下船。
                    \item \zd{这两项义务是强制性的}。
                    \item 免责。共有 17 项，2 项过失（\zd{航行过失、火灾过失}）免责，15 项无过失免责（保险法）。
                    \item \qd{船上自然人是承运人的雇佣人代理人}。船公司远洋运输公司是承运人。
                    \item 便宜！！哈哈哈。火灾责任细分四种。
                          \begin{enumerate}
                            \item 遭雷击不赔，雇佣人代理人过失（如烟头引发）不赔。其他以下都赔。
                            \item \zd{承运人实际过失}导致的。解决实践问题：承运人提供的船舶有问题。
                            \item 私谋（故意纵火）导致火灾。
                          \end{enumerate}
                  \end{enumerate}
                  \tcblower
                  \begin{enumerate}
                    \item 维斯比规则是对海牙规则的非实质性修改。责任限制及最高赔偿限额略高。
                    \item 汉堡规则去掉了过失免责。责任期间：接到交。
                    \item 延迟责任、舱面货、活牲畜、\zd{保函（可考性最强）}。汉堡不允许约定，约定无效。
                    \item 迟延交付的赔偿为迟交货物运费的 2.5 倍，但不应超过应付运费的总额。
                    \item \zd{实际承运人共负连带责任}。
                    \item 诉讼时效，海牙维斯比短 1 年汉堡长 2 年。
                  \end{enumerate}
                \end{kk}
                \begin{kk}{善意保函}
                  \al{神华煤 10W 吨，船老板经验认为只有 8 吨，神华说就有 10W 吨，我给你开保函，你别动笔。这里没有恶意串通，这就是双方发生分歧}。
                  \begin{enumerate}
                    \item \zd{以善意保函换取清洁提单}。
                    \item 在\zd{托运人和承运人}之间有效。\zd{不能对抗收货人不能对抗收货人}。
                    \item 两段法律关系。承收之间的法律文件是提单。托承之间的法律文件才是保函。
                    \item \zd{合同相对性合同相对性合同相对性}。
                  \end{enumerate}
                \end{kk}
        \end{enumerate}
\end{enumerate}

\section{其他国际货物运输法律制度}
\begin{enumerate}
  \item 国际航空货物运输法。考的很少很少。只考过一回。存在最高赔偿限额。
  \item 铁路运输法律制度。一带一路。
        \begin{enumerate}
          \item 运单效力。\zd{不是物权凭证}。
          \item 有单货分离的事实需求，但没有法律基础。要穿越多国，必须随单穿行。通关文牒。
          \item \zd{完全过失连带责任}，无过失才可以免责。
          \item 管货期间为\zd{接到交}期间。
          \item 足额赔偿责任。
        \end{enumerate}
\end{enumerate}
